\section*{Bab \ref{ch:b2:huckel} --- Teori Hückel dan Sistem Terkonjugasi}

\begin{solution}{exo:b2:huckel:1}
Tingkat-tingkat alil ($n = 3$): $\alpha + 1.414\beta$, $\alpha$, $\alpha - 1.414\beta$.
Kation (2 elektron): $2\alpha + 2.828\beta$; radikal (3): $3\alpha + 2.828\beta$;
anion (4): $4\alpha + 2.828\beta$. Elektron-elektron tambahan masuk ke tingkat
nonikatan dan hanya mengubah bagian $\alpha$.
\end{solution}

\begin{solution}{exo:b2:huckel:2}
$2(\alpha + 1.618\beta) + 2(\alpha + 0.618\beta) = 4\alpha + 4.472\beta$.
\end{solution}

\begin{solution}{exo:b2:huckel:3}
Benzena (6), anion siklopentadienil (6), tropilium (6): \omterm{def:b2:huckel:aromatic}{aromatik}. Kation
siklopentadienil (4), siklobutadiena (4), siklooktatetraena datar (8):
\omterm{def:b2:huckel:aromatic}{antiaromatik} jika datar (siklooktatetraena yang sebenarnya berbentuk bak
dan nonaromatik).
\end{solution}

\begin{solution}{exo:b2:huckel:4}
Sebuah segi delapan di dalam lingkaran: tingkat-tingkat $\alpha + 2\beta$,
$\alpha + 1.414\beta$ (dua kali), $\alpha$ (dua kali), $\alpha - 1.414\beta$ (dua
kali), $\alpha - 2\beta$. Delapan elektron: dua pada yang terendah, empat
pada pasangan pertama, satu pada setiap orbital pasangan nonikatan.
\end{solution}

\begin{solution}{exo:b2:huckel:5}
$4.472\beta - 4\beta = 0.472\beta$, yaitu $0.472 \times 75 = \qty{35}{kJ/mol}$,
dibandingkan dengan sekitar \qty{10}{kJ/mol} dari data hidrogenasi: model
Hückel, dengan $\beta$ yang dicocokkan pada benzena, melebih-lebihkan
konjugasi rantai terbuka.
\end{solution}

\begin{solution}{exo:b2:huckel:6}
$p_{12} = p_{34} = 0.894$, $p_{23} = 0.447$: ikatan ujung C1--C2 dan C3--C4
paling pendek, ikatan tengah lebih panjang tetapi lebih pendek daripada
ikatan tunggal.
\end{solution}

\begin{solution}{exo:b2:huckel:7}
Anion, enam elektron: $2 \times 2 + 4 \times 0.618$, $E_\pi = 6\alpha + 6.472\beta$,
kulit tertutup: \omterm{def:b2:huckel:aromatic}{aromatik}. Kation, empat: $2 \times 2 + 2 \times 0.618$,
$4\alpha + 5.236\beta$, dengan satu elektron pada setiap orbital pasangan
terdegenerasi: \omterm{def:b2:huckel:aromatic}{antiaromatik}, sangat tidak stabil.
\end{solution}

\begin{solution}{exo:b2:huckel:8}
Cincin beranggota tujuh: $\alpha + 2\beta$, $\alpha + 1.247\beta$ (dua kali), lalu
tingkat-tingkat antiikatan. Enam elektron, seperti pada kation
\ce{C7H7+}, tepat mengisi tingkat-tingkat ikatan: sebuah kation
\omterm{def:b2:huckel:aromatic}{aromatik}. Senyawa itu mudah terionisasi menjadi tropilium dan bromida.
\end{solution}

\begin{solution}{exo:b2:huckel:9}
$1.802 + 1.247 + 0.445 - 0.445 - 1.247 - 1.802 = 0$: jumlah energinya adalah $6\alpha$.
\end{solution}

\begin{solution}{exo:b2:huckel:10}
$n = 6$: $m_k = 2\cos(k\pi/7) = 1.802$, $1.247$, $0.445$, $-0.445$, $-1.247$,
$-1.802$. Enam elektron:
$E_\pi = 6\alpha + 2(1.802 + 1.247 + 0.445)\beta = 6\alpha + 6.988\beta$;
delokalisasi $0.988\beta$.
\end{solution}

\begin{solution}{exo:b2:huckel:11}
Pada bujur sangkar, kedua elektron berenergi tertinggi masing-masing
menempati satu orbital pasangan nonikatan yang terdegenerasi.
Meregangkan dua ikatan yang berhadapan menghilangkan degenerasi itu:
satu orbital turun, yang lain naik, dan kedua elektron berpasangan pada
yang lebih rendah, yang menurunkan energi. Molekul itu menetap sebagai
persegi panjang dengan dua ikatan rangkap yang terlokalisasi; molekul itu
tetap \omterm{def:b2:huckel:aromatic}{antiaromatik} dan sangat reaktif.
\end{solution}

\begin{solution}{exo:b2:huckel:12}
Matriks (dalam satuan $\beta$, dengan $\alpha$ sebagai titik asal)
$\begin{pmatrix}0 & 1\\ 1 & 1\end{pmatrix}$: $m = 1.618$ dan $-0.618$,
$E = \alpha + 1.618\beta$ (ikatan) dan $\alpha - 0.618\beta$. Untuk \omterm{def:b2:diatomic-mos:bonding}{orbital
ikatan} $c_X = 1.618c_C$, sehingga $c_C^2 = 0.276$, $c_X^2 = 0.724$; dengan dua
elektron, $q_C = 0.55$, $q_X = 1.45$: kerapatan $\pi$ tertarik ke X.
\end{solution}

\begin{solution}{pb:b2:huckel:1}
\textbf{1.} $-123.1 - (-4.3) = \qty{-118.8}{kJ/mol}$.
\textbf{2.} $-123.1 - 104.6 = \qty{-227.7}{kJ/mol}$.
\textbf{3.} $-123.1 - 82.9 = \qty{-206.0}{kJ/mol}$.
\textbf{4.} $3 \times (-118.8) = \qty{-356.4}{kJ/mol}$.
\textbf{5.} $356.4 - 206.0 = \qty{150}{kJ/mol}$.
\textbf{6.} Diena melepaskan $2 \times 118.8 - 227.7 = \qty{10}{kJ/mol}$ lebih
sedikit daripada dua ikatan rangkap terisolasi: konjugasi rantai terbuka
memberikan keuntungan kecil, sedangkan cincin tertutup benzena lima belas
kali lebih besar.
\textbf{7.} Matriks $6 \times 6$ dengan 1 di antara setiap atom dan kedua
tetangganya (1--2, 2--3, \dots, 6--1) dan 0 di tempat lain.
\textbf{8.} $m = 2\cos(2\pi k/6)$: $2$, $1$, $1$, $-1$, $-1$, $-2$: $E = \alpha + 2\beta$,
$\alpha + \beta$ (dua kali), $\alpha - \beta$ (dua kali), $\alpha - 2\beta$.
\textbf{9.} Dua elektron pada $\alpha + 2\beta$, empat pada pasangan
$\alpha + \beta$.
\textbf{10.} $E_\pi = 6\alpha + 8\beta$.
\textbf{11.} $3(2\alpha + 2\beta) = 6\alpha + 6\beta$.
\textbf{12.} $2\beta$, yaitu stabilisasi sebesar $2|\beta|$.
\textbf{13.} $2 + 1 + 1 - 1 - 1 - 2 = 0$: keenam tingkat berjumlah $6\alpha$.
\textbf{14.} $2|\beta| = \qty{150}{kJ/mol}$: $|\beta| = \qty{75}{kJ/mol}$.
\textbf{15.} $0.472 \times 75 = \qty{35}{kJ/mol}$, dibandingkan dengan
\qty{10}{kJ/mol}: model yang dicocokkan pada benzena melebih-lebihkan rantai
terbuka.
\textbf{16.} Orbital-orbital terisi dapat diambil sebagai
$\psi_1 = (1, 1, 1, 1, 1, 1)/\sqrt6$,
$\psi_2 = (2, 1, -1, -2, -1, 1)/\sqrt{12}$, dan $\psi_3 = (0, 1, 1, 0, -1, -1)/2$.
Untuk ikatan 1--2: $p_{12} = 2(\frac16 + \frac{2}{12} + 0) = \frac23$;
berdasarkan simetri setiap ikatan mempunyai $p = 2/3$.
\textbf{17.} Di antara ikatan ujung (0.894) dan ikatan tengah (0.447)
butadiena.
\textbf{18.} Keenam ikatan setara berdasarkan simetri; panjangnya,
\qty{139.7}{pm}, terletak di antara ikatan rangkap etena (\qty{133.9}{pm})
dan ikatan tunggal etana (\qty{153.6}{pm}).
\textbf{19.} $0.988|\beta| = \qty{74}{kJ/mol}$, separuh nilai benzena:
menutup cincin menggandakan \omterm{def:b2:huckel:delocalisation}{energi delokalisasi}.
\textbf{20.} $\alpha + 2\beta$, $\alpha + 1.414\beta$ (dua kali), $\alpha$ (dua kali),
$\alpha - 1.414\beta$ (dua kali), $\alpha - 2\beta$.
\textbf{21.} Kedua elektron terakhir masing-masing masuk ke satu orbital
pasangan nonikatan: kulit terbuka, \omterm{def:b2:huckel:aromatic}{antiaromatik}.
\textbf{22.} Molekul itu melipat menjadi bentuk bak: orbital-orbital p
ikatan-ikatan rangkap yang bertetangga tidak lagi sejajar, ikatan-ikatannya
berselang-seling, dan molekul itu berperilaku sebagai poliena.
\textbf{23.} Empat elektron: pasangan setengah terisi yang sama, lebih buruk
karena molekul itu tidak dapat menghindari kedataran; molekul itu
terdistorsi menjadi persegi panjang dan sangat reaktif.
\textbf{24.} Sistem terkonjugasi monosiklik yang datar bersifat \omterm{def:b2:huckel:aromatic}{aromatik}
dengan $4N + 2$ elektron $\pi$, \omterm{def:b2:huckel:aromatic}{antiaromatik} dengan $4N$.
\textbf{25.} $\boldsymbol{|\beta| \approx \qty{75}{kJ/mol}}$.
\end{solution}
